\documentclass{article}
\usepackage{geometry}
\geometry{margin=1.5cm, vmargin={0pt,1cm}}
\setlength{\topmargin}{-1cm}
\setlength{\headheight}{12.5pt}
\setlength{\paperheight}{29.7cm}
\setlength{\textheight}{25.3cm}

% useful packages.
\usepackage[UTF8]{ctex}
\usepackage{fancyhdr}
\usepackage{layout}
\usepackage{luacode}
\usepackage{listings}
\usepackage{xcolor}

% Configure listings for Python code highlighting
\lstset{
  language=Python,
  basicstyle=\ttfamily\small,
  keywordstyle=\bfseries\color{blue},
  commentstyle=\color{gray},
  stringstyle=\color{red},
  breaklines=true,
  showstringspaces=false,
  tabsize=4,
  frame=single,
  framesep=5pt,
  rulecolor=\color{gray},
  backgroundcolor=\color{white},
  captionpos=b,
  numberstyle=\tiny\color{gray},
  numbers=left,
  stepnumber=5,
}

\lstnewenvironment{plaintext}[1][]{
  \lstset{
    language=none,
    basicstyle=\ttfamily,
    keywordstyle=,
    commentstyle=,
    stringstyle=,
    numbers=none,
    frame=none,
    backgroundcolor=\color{white},
    #1                % 允许用户自行覆盖默认设置
  }
}{}

% import common macros
% % % % % % % % % % % % % % % % % % % % % % % % % % % % % % % %
% Common Macros for LaTeX
% % % % % % % % % % % % % % % % % % % % % % % % % % % % % % % %

% Useful packages
\usepackage{amsfonts}
\usepackage{amsmath}
\usepackage{amssymb}
\usepackage{amsthm}
\usepackage{enumerate}
\usepackage{graphicx}
\usepackage{multicol}
\usepackage[ruled,linesnumbered]{algorithm2e}

% Math operators and common symbols
\newcommand{\dif}{\mathrm{d}}
\newcommand{\innerProd}[1]{\left\langle #1 \right\rangle}
\newcommand{\difFrac}[2]{\frac{\dif #1}{\dif #2}}
\newcommand{\pdfFrac}[2]{\frac{\partial #1}{\partial #2}}
\newcommand{\T}{\mathrm{T}}
\newcommand{\norm}[1]{\left\| #1 \right\|}
\newcommand{\abs}[1]{\left| #1 \right|}

% Vector and Matrix shortcuts
\newcommand{\vecTwo}[2]{
  \begin{bmatrix}
    #1 \\
    #2
\end{bmatrix}}
\newcommand{\vecThree}[3]{
  \begin{bmatrix}
    #1 \\
    #2 \\
    #3
\end{bmatrix}}
\newcommand{\vecTwoH}[2]{
  \begin{bmatrix}
    #1 & #2
  \end{bmatrix}
}
\newcommand{\vecThreeH}[3]{
  \begin{bmatrix}
    #1 & #2 & #3
  \end{bmatrix}
}

% Common fields
\newcommand{\R}{\mathbb{R}}
\newcommand{\C}{\mathbb{C}}
\newcommand{\Q}{\mathbb{Q}}
\newcommand{\Z}{\mathbb{Z}}
\newcommand{\N}{\mathbb{N}}

% Common Operators
\renewcommand{\Re}{\mathrm{Re}}
\renewcommand{\Im}{\mathrm{Im}}

% Theorem-like environments
\theoremstyle{definition}
\newtheorem{theorem}{Theorem}[section]
\newtheorem{lemma}[theorem]{Lemma}
\newtheorem{corollary}[theorem]{Corollary}
\newtheorem{definition}[theorem]{Definition}
\newtheorem{remark}[theorem]{Remark}
\newtheorem{exercise}{Exercise}[section]

% Support for individual Chinese characters using fontspec (for LuaLaTeX)
\usepackage{fontspec}
\newfontfamily\zhfont{SimSun} % Search system fonts by name
\newcommand{\zhstr}[1]{{\zhfont #1}}

% Solutions and proofs
\AtBeginDocument{
  \ifdefined\ctexset
  \newenvironment{solution}{
  \begin{proof}[解]}{
  \end{proof}}
  \else
  \newenvironment{solution}{
  \begin{proof}[Solution]}{
  \end{proof}}
  \fi
}

\usepackage{luacode}
% Utility to get environment variables (requires lualatex)
\newcommand{\getEnv}[2]{\directlua{tex.print(os.getenv("#1") or "#2")}}


% Name and uID
\newcommand{\name}{\getEnv{NAME}{NAME}}
\newcommand{\uid}{\getEnv{UID}{UID}}
\newcommand{\course}{数值代数}
\newcommand{\hwNum}{2}

\begin{document}

\pagestyle{fancy}
\fancyhead{}
\lhead{\name (\uid)}
\chead{\course \#\hwNum}
\rhead{\today}

\section{题目 1}

\subsection*{(1) 上（下）三角阵的逆是上（下）三角阵}

\begin{proof}
  设 $A$ 是 $n \times n$ 上三角阵，$A$ 可逆。则 $\det(A) \neq 0$，即 $A$ 的对角线元素都非零。

  设 $A^{-1} = B = (b_{ij})$，由 $AB = I$ 得：
  \[
    \sum_{k=1}^{n} a_{ik} b_{kj} = \delta_{ij}
  \]

  对于 $i > j$，考虑第 $i$ 行：若 $A$ 是上三角阵，则 $a_{ik} = 0$（$k < i$）。因此：
  \[
    \sum_{k=i}^{n} a_{ik} b_{kj} = \delta_{ij}
  \]

  当 $i > j$ 时 $\delta_{ij} = 0$，所以：
  \[
    \sum_{k=i}^{n} a_{ik} b_{kj} = 0
  \]

  用逆向归纳法，对 $k = n, n-1, \ldots, i$：
  \[
    a_{ik} b_{kj} = 0 \quad (k \geq i)
  \]

  由于 $a_{ii} \neq 0$，得 $b_{ij} = 0$（$i > j$），故 $A^{-1}$ 是上三角阵。

  下三角阵的情况类似可证。
\end{proof}

\subsection*{(2) 单位三角阵的逆是单位三角阵}

\begin{proof}
  设 $L$ 是单位下三角阵（对角线元素为 1），则 $L$ 是下三角阵，由 (1) 知 $L^{-1}$ 也是下三角阵。

  设 $L^{-1} = M = (m_{ij})$，由 $LM = I$ 得：
  \[
    \sum_{k=1}^{n} l_{ik} m_{kj} = \delta_{ij}
  \]

  当 $i = j$ 时：
  \[
    l_{ii} m_{ii} + \sum_{k=i+1}^{n} l_{ik} m_{kj} = 1
  \]

  由于 $l_{ii} = 1$，$l_{ik} = 0$（$k > i$），得：
  \[
    m_{ii} = 1
  \]

  因此 $L^{-1}$ 的对角线元素也都为 1，即 $L^{-1}$ 是单位下三角阵。

  上三角阵情况类似。
\end{proof}

\section{题目 2}

\subsection*{算法}

请你设计一个变形的高斯消元法，从右往左消去 $A$ 的列得到三角分解 $A=UL$。
\begin{solution}

  矩阵 $U$ 从单位矩阵开始，每次消去时更新第 $j$ 列的上三角部分；矩阵 $L$ 从 $A$ 开始，每次消去
  时将第 $j$ 列的上三角部分消去掉。最终 $U$ 是单位上三角阵，$L$ 是下三角阵，并且满足 $A = UL$。

  \begin{algorithm}[H]
    \DontPrintSemicolon
    \caption{从右向左的列消去，得到 $A=UL$}
    \KwIn{非奇异矩阵 $A \in \mathbb{R}^{n \times n}$}
    \KwOut{单位上三角阵 $U$ 与下三角阵 $L$，满足 $A=UL$}
    $U = I_n,\ L = A$\;
    \For{$j = n : -1 :2$}{
      $U(1:j-1, j) = L(1:j-1, j) / L(j, j)$\;
      $L(1:j-1, :) = L(1:j-1, :) - U(1:j-1, j)  L(j, :)$\;
    }
    \Return{$U,L$}\;
  \end{algorithm}

\end{solution}

\subsection*{唯一性证明}

\begin{proof}
  假设存在两个这样的分解：
  \[
    A = U_1 L_1 = U_2 L_2
  \]

  其中 $U_1, U_2$ 为单位上三角阵，$L_1, L_2$ 为下三角阵。

  则：
  \[
    U_1 L_1 = U_2 L_2 \implies U_2^{-1} U_1 = L_2 L_1^{-1}
  \]

  等式左边 $U_2^{-1} U_1$ 是单位上三角阵的乘积，仍为单位上三角阵。

  等式右边 $L_2 L_1^{-1}$ 是下三角阵的乘积，由题目1知 $L_1^{-1}$ 也是下三角阵，故 $L_2 L_1^{-1}$ 为下三角阵。

  一个矩阵既是单位上三角阵又是下三角阵，只能是单位矩阵 $I$。

  因此：
  \[
    U_2^{-1} U_1 = I \implies U_1 = U_2
  \]
  \[
    L_2 L_1^{-1} = I \implies L_1 = L_2
  \]

  故分解唯一。
\end{proof}

\section{题目 3}

\subsection*{题面}

设
\[
  A =
  \begin{pmatrix}1&4&7\\2&5&8\\3&6&10
  \end{pmatrix}, \quad
  b =
  \begin{pmatrix}1\\1\\1
  \end{pmatrix}
\]

\subsection*{(1) $A$ 的 $LU$ 分解}
\begin{solution}

  对矩阵 $A$ 进行不选主元的高斯消元法。

  \textbf{第一步：消去第 1 列}

  选择 $a_{11} = 1$ 作为主元：
  \begin{align*}
    m_{21} &= \frac{2}{1} = 2, & m_{31} &= \frac{3}{1} = 3
  \end{align*}
  执行行消元：
  \begin{align*}
    R_2 &:= R_2 - 2R_1 \implies (0, -3, -6) \\
    R_3 &:= R_3 - 3R_1 \implies (0, -6, -11)
  \end{align*}
  得到矩阵：$
  \begin{pmatrix}1&4&7\\0&-3&-6\\0&-6&-11
  \end{pmatrix}$

  \textbf{第二步：消去第 2 列}

  选择 $a_{22} = -3$ 作为主元：
  \[
    m_{32} = \frac{-6}{-3} = 2
  \]
  执行行消元：
  \[
    R_3 := R_3 - 2R_2 \implies (0, 0, 1)
  \]

  \textbf{LU 分解结果：}
  \[
    \boxed{
      L =
      \begin{pmatrix}1&0&0\\2&1&0\\3&2&1
      \end{pmatrix}, \quad
      U =
      \begin{pmatrix}1&4&7\\0&-3&-6\\0&0&1
      \end{pmatrix}
    }
  \]

  \textbf{求解 $Ax = b$：}

  首先解 $Ly = b$：
  \begin{align*}
    y_1 &= 1 \\
    2y_1 + y_2 &= 1 \implies y_2 = -1 \\
    3y_1 + 2y_2 + y_3 &= 1 \implies y_3 = 0
  \end{align*}

  再解 $Ux = y$：
  \begin{align*}
    x_3 &= 0 \\
    -3x_2 - 6x_3 &= -1 \implies x_2 = \frac{1}{3} \\
    x_1 + 4x_2 + 7x_3 &= 1 \implies x_1 = -\frac{1}{3}
  \end{align*}

  \textbf{解：} $\boxed{x = \left(-\dfrac{1}{3}, \dfrac{1}{3}, 0\right)^T}$
\end{solution}

\subsection*{(2) 列主元分解 $PA = LU$}
\begin{solution}

  采用列主元高斯消元法。

  \textbf{第一步：选主元}

  第 1 列的元素：$|1| = 1, |2| = 2, |3| = 3$，最大为 3 在第 3 行。

  进行行交换 $P_1$（交换第 1 和 3 行）：
  \[
    P_1 A =
    \begin{pmatrix}3&6&10\\2&5&8\\1&4&7
    \end{pmatrix}
  \]

  消元，乘子 $m_{21} = \frac{2}{3}, m_{31} = \frac{1}{3}$：
  \[
    \begin{pmatrix}3&6&10\\0&1&\frac{4}{3}\\0&2&\frac{11}{3}
    \end{pmatrix}
  \]

  \textbf{第二步：选主元}

  第 2 列第 2 行以下的元素：$|1| = 1, |2| = 2$，最大为 2 在第 3 行。

  进行行交换 $P_2$（交换第 2 和 3 行）：
  \[
    \begin{pmatrix}3&6&10\\0&2&\frac{11}{3}\\0&1&\frac{4}{3}
    \end{pmatrix}
  \]

  消元，乘子 $m_{32} = \frac{1}{2}$：
  \[
    U =
    \begin{pmatrix}3&6&10\\0&2&\frac{11}{3}\\0&0&-\frac{1}{2}
    \end{pmatrix}
  \]

  \textbf{列主元分解结果：}

  排列矩阵 $P$（对应行变换序列）：
  初始排列为 $[1,2,3]$，第一次交换后为 $[3,2,1]$，第二次交换后为 $[3,1,2]$

  \[
    \boxed{
      P =
      \begin{pmatrix}
        0 & 0 & 1 \\
        1 & 0 & 0 \\
        0 & 1 & 0
      \end{pmatrix}, \quad
      L =
      \begin{pmatrix}
        1 & 0 & 0 \\
        \frac{1}{3} & 1 & 0 \\
        \frac{2}{3} & \frac{1}{2} & 1
      \end{pmatrix}, \quad
      U =
      \begin{pmatrix}
        3 & 6 & 10 \\
        0 & 2 &\frac{11}{3} \\
        0 & 0 & -\frac{1}{2}
      \end{pmatrix}
    }
  \]

  \textbf{求解 $Ax = b$：}

  $Pb =
  \begin{pmatrix}1\\1\\1
  \end{pmatrix}$（因为 $b$ 的三个分量相等）

  解 $Ly = Pb$：
  \begin{align*}
    y_1 &= 1 \\
    \frac{1}{3} + y_2 &= 1 \implies y_2 = \frac{2}{3} \\
    \frac{2}{3} + \frac{1}{2} \cdot \frac{2}{3} + y_3 &= 1 \implies y_3 = 0
  \end{align*}

  解 $Ux = y$：
  \begin{align*}
    -\frac{1}{2}x_3 &= 0 \implies x_3 = 0 \\
    2x_2 + \frac{11}{3} \cdot 0 &= \frac{2}{3} \implies x_2 = \frac{1}{3} \\
    3x_1 + 6 \cdot \frac{1}{3} + 10 \cdot 0 &= 1 \implies x_1 = -\frac{1}{3}
  \end{align*}

  \textbf{解：} $\boxed{x = \left(-\dfrac{1}{3}, \dfrac{1}{3}, 0\right)^T}$
\end{solution}

\subsection*{(3) 全主元分解 $PAQ = LU$}

\begin{solution}

  全主元高斯消元法在每一步都选整个右下子矩阵中绝对值最大的元素。

  \textbf{第一步：选主元}

  全矩阵中最大元素为 10，在位置 $(3,3)$。

  行交换：交换第 1 和 3 行；列交换：交换第 1 和 3 列

  \[
    P_1 A Q_1 =
    \begin{pmatrix}
      10 & 6 & 3 \\
      8 & 5 & 2 \\
      7 & 4 & 1
    \end{pmatrix}
  \]

  消元，乘子 $m_{21} = \frac{4}{5}, m_{31} = \frac{7}{10}$：
  \[
    \begin{pmatrix}
      10 & 6 & 3 \\
      0 & \frac{1}{5} & -\frac{2}{5} \\
      0 & -\frac{1}{5} & -\frac{11}{10}
    \end{pmatrix}
  \]

  \textbf{第二步：选主元}

  右下角 $2 \times 2$ 子矩阵中最大元素为 $|-\frac{11}{10}| = \frac{11}{10}$ 在位置 $(3,3)$。

  交换第 2 和第 3 行，交换第 2 和第 3 列：

  \[
    \begin{pmatrix}
      10 & 3 & 6 \\
      0 & -\frac{11}{10} & -\frac{1}{5} \\
      0 & -\frac{2}{5} & \frac{1}{5}
    \end{pmatrix}
  \]

  消元，乘子 $m_{32} = \frac{4}{11}$：
  \[
    U =
    \begin{pmatrix}
      10 & 3 & 6 \\
      0 & -\frac{11}{10} & -\frac{1}{5} \\
      0 & 0 & \frac{3}{11}
    \end{pmatrix}
  \]

  \textbf{全主元分解结果：}

  列排列矩阵（第 1 和 3 列交换后，第 2 和 3 列交换）：
  \[
    Q =
    \begin{pmatrix}
      0 & 1 & 0 \\
      0 & 0 & 1 \\
      1 & 0 & 0
    \end{pmatrix}
  \]

  行排列矩阵（第 1 和 3 行交换后，第 2 和 3 行交换）：
  \[
    P =
    \begin{pmatrix}
      0 & 0 & 1 \\
      1 & 0 & 0 \\
      0 & 1 & 0
    \end{pmatrix}
  \]

  下三角矩阵：
  \[
    L =
    \begin{pmatrix}
      1 & 0 & 0 \\
      \frac{7}{10} & 1 & 0 \\
      \frac{4}{5} & \frac{4}{11} & 1
    \end{pmatrix}
  \]

  于是：
  \[
    \boxed{
      P =
      \begin{pmatrix}
        0 & 0 & 1 \\
        1 & 0 & 0 \\
        0 & 1 & 0
      \end{pmatrix}, \quad
      Q =
      \begin{pmatrix}
        0 & 1 & 0 \\
        0 & 0 & 1 \\
        1 & 0 & 0
      \end{pmatrix}, \quad
      L =
      \begin{pmatrix}
        1 & 0 & 0 \\
        \frac{7}{10} & 1 & 0 \\
        \frac{4}{5} & \frac{4}{11} & 1
      \end{pmatrix}, \quad
      U =
      \begin{pmatrix}
        10 & 3 & 6 \\
        0 & -\frac{11}{10} & -\frac{1}{5} \\
        0 & 0 & \frac{3}{11}
      \end{pmatrix}
    }
  \]

  \textbf{求解 $Ax = b$：}

  由 $PAQz = Pb$ 和 $x = Qz$，先求 $z$。

  $Pb =
  \begin{pmatrix}1\\1\\1
  \end{pmatrix}$

  解 $Ly = Pb$：
  \begin{align*}
    y_1 &= 1 \\
    \frac{7}{10} + y_2 &= 1 \implies y_2 = \frac{3}{10} \\
    \frac{4}{5} + \frac{4}{11} \cdot \frac{3}{10} + y_3 &= 1 \implies
    y_3 = \frac{1}{11}
  \end{align*}

  解 $Uz = y$：
  \begin{align*}
    \frac{3}{11} z_3 &= \frac{1}{11} \implies z_3 = \frac{1}{3} \\
    -\frac{11}{10} z_2 - \frac{1}{5} \cdot \frac{1}{3} &=
    \frac{3}{10} \implies z_2 = -\frac{1}{3} \\
    10 z_1 + 3 \cdot \left(-\frac{1}{3}\right) + 6 \cdot \frac{1}{3} &= 1
    \implies z_1 = 0
  \end{align*}

  通过列排列恢复解：
  \[
    x = Qz = \boxed{\left( -\frac{1}{3}, \frac{1}{3}, 0\right)^T}
  \]
\end{solution}

\section{上机习题 1.1}

教材 39 页第一题：实现不选主元和列主元高斯消元法并求解方程。这里使用 Python 语言实现，并与 SciPy 库中的求解函数进行比较。

\lstinputlisting[language=Python,caption=高斯消元法实现（不选主元和列主元）]{exercise1.py}

\begin{lstlisting}[language={},caption=运行结果]
Time taken without pivoting: 0.001011 seconds
Time taken with column pivoting: 0.001655 seconds
Time taken using SciPy function: 0.284788 seconds
Difference for no pivoting: 5.960464477539076e-08
Difference for column pivoting: 2.5121479338940403e-15
Difference for SciPy's solution: 1.7763568394002505e-15
\end{lstlisting}

这里容易看到，列主元高斯消元法的结果与 SciPy 的求解函数非常接近（误差在机器精度范围内），而不选主元的结果误差较大，说明了选主元的重要性。
事实上，这里不选主元消元法精度会发生爆炸，从第 70 行左右开始误差达到 $10^8$ 的量级，导致最终结果误差达到 $10^{-8}$ 的量级，而列主元消元法的误差保持在机器精度范围内。

\end{document}
